\scanpage{64}
\mspage{70}

\margin{\figure{\begin{tikzpicture}[x=0.24cm,y=0.24cm]
\draw[gray,line width=0.7pt,line cap=round] (0.62,1.9) .. controls (0.86,1.61) and (0.84,1.36) .. (0.9,1.2) .. controls (1.28,1.51) and (1.32,1.64) .. (1.39,1.74);
\draw[gray,line width=0.7pt,line cap=round] (0.05,1.57) .. controls (0.51,1.3) and (0.92,1.18) .. (1.4,1.02);
\draw[gray,line width=0.7pt,line cap=round] (0.91,1.22) .. controls (0.98,0.74) and (0.98,0.12) .. (0.53,0.13) .. controls (0.01,0.15) and (0.33,0.75) .. (0.91,1.2);
\end{tikzpicture}}}

\deleted{Appliquons ceci au foncteur déduit de $f:X\to Y$ par le changement de base $Y_{/y_0}\longrightarrow Y$, \unsure{d'où}} \inserted{Considérons} le diagramme à carrés cartésiens

\diagram{\begin{tikzcd}[ampersand replacement=\&]
X_{y_0}\arrow[r,"q"]\arrow[d,"p"] \& X_{/y_0}\arrow[r]\arrow[d,"f_0"] \& X\arrow[d,"f"]\\
\{y_0\}\arrow[r] \& Y_{/y_0}\arrow[r] \& Y,
\end{tikzcd}}

\deleted{de cette façon on trouve} \inserted{L'hyp. c$''$) prouve} que pour tout $F\in\unsure{\mathbb D|\underline D}(X_{/y_0})$ l'on a $H^\bullet_{\mathbb D}(X_{/y_0},F)\xrightarrow{\sim}H^\bullet_{\mathbb D}(X_{y_0},F)$. Par le th. 1, \unsure{cela|celà} implique que $q$ est une $\mathbb D$-équiv. co-stricte.

\uline{Cor. 1} Énoncé dual du th., pour les \unsure{deux|de.} foncteurs $X_{y_0}\longrightarrow{}_{y_0\backslash}X$, $X_{/x_0}\longrightarrow Y_{/y_0}$, et le chang\unsure{emt|ment} de \deleted{\unsure{base}} dans $f_!(F)$ \ldots

\uline{Déf. 4} Sous les conditions \inserted{équiv.} du th. 1, on dit que $f$ est \deleted{\unsure{cohomologiquemt|catérologiquemt}} \uline{$\mathbb D$-propre}. La condition duale s'appelle $f$ $\mathbb D$-lisse.

\deleted{\uline{Prop.} Notons d'abord que chgt de base, par composition}

\uline{Cor. 2} : Soit $f:X\longrightarrow Y$ \unsure{dans|du} $\mathbb D$. Si $f$ est $\mathbb D$-propre \deleted{\unsure{fibré}} ou $\mathbb D$-lisse, et à fibres $\mathbb D$-contractiles, alors $f$ est une $\mathbb D$-équce stricte et costricte. Mieux\unsure{:|,} $f$ est une $\mathbb D$-équce \uline{universelle} i.e. \deleted{\unsure{pour}} c'est une

% slanted
\margin{\uline{Cor 2$'$} \uline{\unsure{généralisé|généraliser}} : Soit \diagram{\begin{tikzcd}[ampersand replacement=\&,column sep=tiny,row sep=tiny]
X\arrow[rr,"f"]\arrow[dr] \& {} \& Y\arrow[dl]\\
{} \& S \& {}
\end{tikzcd}} avec $X,Y$ \unsure{propres|fibrés} sur $S$. Si $f$ induit une $\mathbb D$-équiv. sur les fibres, \unsure{c'est} une $\mathbb D$-équiv. \unsure{Idem} dans le cas lisse. \unsure{C'est le principe|C'est le privilège} \unsure{des conditions|des conclusions}.}

\deleted{Il suffit de voir que $f$ est une $\mathbb D$-équivalence, cela le restera par tt chgt de base \unsure{vu que} les énoncés \unsure{sur} $\mathbb D$-propre et $\mathbb D$-lisse.} $\mathbb D$-équiv, et le reste par \uline{tout} chgt de \inserted{base $Y'\to Y$.}
% archive sheet number: 63
